\documentclass[lualatex,ja=standard,a4paper,12pt]{bxjsarticle}
\usepackage{amsmath, amssymb, amsthm}
\usepackage{geometry}
\usepackage{xcolor}
\usepackage{tcolorbox}
\usepackage{booktabs}
\usepackage{array}
\usepackage{pgfplots}
\pgfplotsset{compat=1.18}
\pgfmathdeclarefunction{gauss}{2}{%
  \pgfmathparse{1/(#2*sqrt(2*pi))*exp(-((x-#1)^2)/(2*#2^2))}%
}

\geometry{margin=25mm}

\tcbuselibrary{theorems, skins, breakable}

% ---- カラーボックス定義 ----
\newtcolorbox{defbox}[1]{
  colback=blue!5, colframe=blue!60!black,
  fonttitle=\bfseries, title={#1}, arc=3pt
}
\newtcolorbox{thmbox}[1]{
  colback=green!5, colframe=green!50!black,
  fonttitle=\bfseries, title={#1}, arc=3pt
}
\newtcolorbox{pointbox}{
  colback=orange!8, colframe=orange!70!black,
  fonttitle=\bfseries, title={ポイント}, arc=3pt
}
\newtcolorbox{exbox}[1]{
  colback=gray!6, colframe=gray!50,
  fonttitle=\bfseries, title={例：#1}, arc=3pt,
  breakable
}

\title{{\large Hoff 第5章} \\[6pt]
       \textbf{5.4\; Bias, Variance and Mean Squared Error} \\[4pt]
       {\normalsize ゼミ資料}}
\author{}
\date{}

\begin{document}
\maketitle
\thispagestyle{empty}

% ===========================
\section{この節のモチベーション}
% ===========================

前節までで，正規モデルにおける$\theta$の事後平均（ベイズ推定量）

\[
  \hat{\theta}_b = \frac{\kappa_0}{\kappa_0+n}\,\mu_0
                 + \frac{n}{\kappa_0+n}\,\bar{y}
                 = w\bar{y} + (1-w)\mu_0,
  \qquad w = \frac{n}{\kappa_0+n}
\]

を導いた．これは事前平均$\mu_0$と標本平均$\bar{y}$の加重平均になっている．

\medskip
\noindent
\textbf{素朴な疑問：} 推定量として「標本平均 $\hat{\theta}_e = \bar{y}$」と
「ベイズ推定量 $\hat{\theta}_b$」のどちらが良いのか？

\medskip
\noindent
この問いに答えるために，\textbf{バイアス・分散・平均二乗誤差（MSE）}という枠組みを使う．

% ===========================
\section{推定量の標本特性}
% ===========================

\begin{defbox}{点推定量（Point Estimator）}
未知パラメータ$\theta$の\textbf{点推定量}とは，データ$(y_1,\ldots,y_n)$を
パラメータ空間$\Theta$の一点に変換する関数のことである．
\end{defbox}

\medskip
推定量はデータの関数なので，\textbf{それ自体が確率変数}である．
データを無限回集め直したとしたら，推定値はどのように分布するか？
---これが\textbf{標本特性（sampling properties）}の考え方である．

\bigskip
今回比較する2つの推定量：

\[
  \hat{\theta}_e = \bar{y}
  \qquad \text{（標本平均）}
\]
\[
  \hat{\theta}_b = w\bar{y} + (1-w)\mu_0
  \qquad \text{（ベイズ推定量：事後平均）}
\]

真の母平均を$\theta_0$，母分散を$\sigma^2$（既知と仮定）とおく．

% ===========================
\section{バイアス（偏り）}
% ===========================

\begin{defbox}{バイアス（Bias）}
推定量$\hat{\theta}$のバイアスを次で定義する：
\[
  \mathrm{Bias}[\hat{\theta} \mid \theta_0]
  = E[\hat{\theta} \mid \theta_0] - \theta_0.
\]
バイアスが0の推定量を\textbf{不偏推定量（unbiased estimator）}と呼ぶ．
\end{defbox}

\subsection*{各推定量のバイアスを計算する}

\paragraph{標本平均 $\hat{\theta}_e$ のバイアス}

\[
  E[\hat{\theta}_e \mid \theta_0]
  = E\!\left[\frac{1}{n}\sum_{i=1}^n y_i \;\middle|\; \theta_0\right]
  = \frac{1}{n}\sum_{i=1}^n \theta_0 = \theta_0
\]

\[
  \therefore\quad \mathrm{Bias}[\hat{\theta}_e \mid \theta_0] = 0
  \qquad \Rightarrow \quad \hat{\theta}_e \text{ は\textbf{不偏}}
\]

\paragraph{ベイズ推定量 $\hat{\theta}_b$ のバイアス}

\[
  E[\hat{\theta}_b \mid \theta_0]
  = w\,E[\bar{y} \mid \theta_0] + (1-w)\mu_0
  = w\theta_0 + (1-w)\mu_0
\]

\[
  \therefore\quad \mathrm{Bias}[\hat{\theta}_b \mid \theta_0]
  = w\theta_0 + (1-w)\mu_0 - \theta_0
  = (1-w)(\mu_0 - \theta_0)
\]

$\mu_0 \neq \theta_0$ のとき，$\hat{\theta}_b$ は\textbf{偏りあり（biased）}．

\begin{pointbox}
不偏性は一見良さそうだが，それだけで推定量の良さは語れない．
たとえば$y_1$だけも$\theta_0$の不偏推定量だが，$\bar{y}$より真値から遠い．
「バイアスがゼロ＝良い推定量」ではない！
\end{pointbox}

% ===========================
\section{平均二乗誤差（MSE）}
% ===========================

推定量 $\hat{\theta}$ が真値 $\theta_0$ にどれだけ近いかを総合的に評価するために，
\textbf{平均二乗誤差（Mean Squared Error; MSE）}を用いる．

\begin{defbox}{MSE（Mean Squared Error）}
\[
  \mathrm{MSE}[\hat{\theta} \mid \theta_0]
  = E\!\left[(\hat{\theta} - \theta_0)^2 \;\middle|\; \theta_0\right]
\]
\end{defbox}

\subsection*{MSE の分解}

$m = E[\hat{\theta} \mid \theta_0]$ とおいて式を展開する：

\begin{align*}
  \mathrm{MSE}[\hat{\theta} \mid \theta_0]
  &= E[(\hat{\theta} - \theta_0)^2 \mid \theta_0] \\
  &= E\!\bigl[(\hat{\theta} - m + m - \theta_0)^2 \mid \theta_0\bigr] \\
  &= E[(\hat{\theta} - m)^2 \mid \theta_0]
   + 2\,\underbrace{E[(\hat{\theta}-m) \mid \theta_0]}_{=\,0}\,(m-\theta_0)
   + (m - \theta_0)^2.
\end{align*}

$E[\hat{\theta} - m \mid \theta_0] = 0$ より第2項は消えるので，

\begin{thmbox}{MSE の分解定理}
\[
  \mathrm{MSE}[\hat{\theta} \mid \theta_0]
  \;=\;
  \underbrace{\mathrm{Var}[\hat{\theta} \mid \theta_0]}_{\text{分散}}
  \;+\;
  \underbrace{\mathrm{Bias}^2[\hat{\theta} \mid \theta_0]}_{\text{バイアスの2乗}}
\]
\end{thmbox}

\begin{pointbox}
MSE は「分散」と「バイアスの2乗」のトレードオフで決まる．\\
バイアスがあっても，分散を十分小さくできれば MSE を改善できる可能性がある．
\end{pointbox}

% ===========================
\section{2つの推定量の MSE 比較}
% ===========================

\subsection*{分散の比較}

\[
  \mathrm{Var}[\hat{\theta}_e \mid \theta_0, \sigma^2]
  = \mathrm{Var}\!\left[\bar{y}\right] = \frac{\sigma^2}{n}
\]

\[
  \mathrm{Var}[\hat{\theta}_b \mid \theta_0, \sigma^2]
  = \mathrm{Var}[w\bar{y}] = w^2\,\frac{\sigma^2}{n}
  < \frac{\sigma^2}{n}
\]

$w < 1$ なので，\textbf{ベイズ推定量の方が分散は常に小さい}．

\subsection*{MSE の比較}

\[
  \mathrm{MSE}[\hat{\theta}_e \mid \theta_0]
  = \frac{\sigma^2}{n} + 0
  = \frac{\sigma^2}{n}
\]

\[
  \mathrm{MSE}[\hat{\theta}_b \mid \theta_0]
  = w^2\,\frac{\sigma^2}{n} + (1-w)^2(\mu_0 - \theta_0)^2
\]

\subsection*{どちらが良いか？}

代数計算により，$\mathrm{MSE}[\hat{\theta}_b] < \mathrm{MSE}[\hat{\theta}_e]$ となる条件は

\[
  \boxed{
    (\mu_0 - \theta_0)^2
    < \sigma^2\!\left(\frac{1}{n} + \frac{2}{\kappa_0}\right)
  }
\]

\paragraph{導出の概略}

$\mathrm{MSE}[\hat{\theta}_b] < \mathrm{MSE}[\hat{\theta}_e]$ を整理すると

\begin{align*}
  w^2\,\frac{\sigma^2}{n} + (1-w)^2(\mu_0-\theta_0)^2
  &< \frac{\sigma^2}{n} \\
  (1-w)^2(\mu_0-\theta_0)^2
  &< (1-w^2)\,\frac{\sigma^2}{n}
   = (1-w)(1+w)\,\frac{\sigma^2}{n} \\
  (1-w)(\mu_0-\theta_0)^2
  &< (1+w)\,\frac{\sigma^2}{n}
\end{align*}

$w = n/(\kappa_0+n)$ を代入すると，
\[
  \frac{1+w}{1-w}
  = \frac{\,1 + \dfrac{n}{\kappa_0+n}\,}{1 - \dfrac{n}{\kappa_0+n}}
  = \frac{\kappa_0+2n}{\kappa_0},
\]
よって
\[
  (\mu_0-\theta_0)^2
  < \frac{\sigma^2}{n} \cdot \frac{\kappa_0+2n}{\kappa_0}
  = \sigma^2\!\left(\frac{1}{n} + \frac{2}{\kappa_0}\right).
\]

\begin{pointbox}
\textbf{解釈：}
事前平均$\mu_0$が真値$\theta_0$から「$\sigma^2(1/n + 2/\kappa_0)$以内に近い」なら，
ベイズ推定量の方が MSE の意味で優れる．

\medskip
事前情報が多少でもあれば（つまり$\mu_0$が大きく外れていなければ），
ベイズ推定量は標本平均よりも良い推定を与えることができる．
\end{pointbox}

% ===========================
\section{例：IQスコアによる比較}
% ===========================

\begin{exbox}{IQスコア}

IQテストは一般母集団で$\mu=100$，$\sigma=15$になるよう設計されている．
ある町の住民から$n$人をサンプリングして，
その町の平均IQスコア$\theta$を推定する．

\medskip
\noindent\textbf{設定：}
\begin{itemize}
  \item 事前分布：$\mu_0 = 100$（一般母集団の平均を使う）
  \item 真のパラメータ（未知）：$\theta_0 = 112$，$\sigma^2 = 169$
  \item 実はこの町は非常に優秀で，真の平均が100よりずっと高い！
\end{itemize}

\medskip
\noindent\textbf{各推定量の MSE：}

\[
  \mathrm{MSE}[\hat{\theta}_e \mid \theta_0] = \frac{169}{n}
\]

\[
  \mathrm{MSE}[\hat{\theta}_b \mid \theta_0]
  = w^2\,\frac{169}{n} + (1-w)^2 \times (112-100)^2
  = w^2\,\frac{169}{n} + (1-w)^2 \times 144
\]

\medskip
\noindent\textbf{図1：相対MSE（$n$の関数として）} \\[4pt]
$w = n/(\kappa_0+n)$ とおくと，相対MSEは
\[
  \frac{\mathrm{MSE}[\hat{\theta}_b]}{\mathrm{MSE}[\hat{\theta}_e]}
  = w^2 + (1-w)^2 \cdot \frac{144\,n}{169}
\]

\begin{center}
\begin{tikzpicture}
\begin{axis}[
  width=8.5cm, height=6.5cm,
  xlabel={サンプルサイズ $n$},
  ylabel={相対MSE},
  xmin=0, xmax=50, ymin=0, ymax=1.15,
  legend pos=south east,
  legend style={font=\small},
  grid=major, grid style={dotted, gray!50},
  tick label style={font=\small},
  label style={font=\small},
]
% κ₀=0: w=1, ratio=1 (常に1)
\addplot[black, dashed, thick, domain=0.5:50, samples=2] {1};
\addlegendentry{$\kappa_0=0$（標本平均）}

% κ₀=1
\addplot[blue!80!black, thick, domain=0.5:50, samples=200]
  { (x/(1+x))^2 + (1/(1+x))^2 * 144*x/169 };
\addlegendentry{$\kappa_0=1$}

% κ₀=2
\addplot[red!80!black, thick, domain=0.5:50, samples=200]
  { (x/(2+x))^2 + (2/(2+x))^2 * 144*x/169 };
\addlegendentry{$\kappa_0=2$}

% κ₀=3
\addplot[green!60!black, thick, domain=0.5:50, samples=200]
  { (x/(3+x))^2 + (3/(3+x))^2 * 144*x/169 };
\addlegendentry{$\kappa_0=3$}
\end{axis}
\end{tikzpicture}
\end{center}

\medskip
\noindent\textbf{図2：推定量の標本分布（$n=5$のとき）} \\[4pt]
$n=5$, $\sigma^2=169$, $\mu_0=100$, $\theta_0=112$ として，各推定量は以下の正規分布に従う：

\smallskip
\begin{center}
\small
\begin{tabular}{lll}
\toprule
推定量 & 期待値 & 標準偏差 \\
\midrule
$\hat\theta_e$ & $112$ & $\sqrt{169/5} \approx 5.81$ \\
$\hat\theta_b\ (\kappa_0=1)$ & $110.0$ & $\approx 4.84$ \\
$\hat\theta_b\ (\kappa_0=2)$ & $108.6$ & $\approx 4.15$ \\
$\hat\theta_b\ (\kappa_0=3)$ & $107.5$ & $\approx 3.63$ \\
\bottomrule
\end{tabular}
\end{center}

\begin{center}
\begin{tikzpicture}
\begin{axis}[
  width=8.5cm, height=6.5cm,
  xlabel={IQスコアの推定値},
  ylabel={密度},
  xmin=90, xmax=135,
  ymin=0, ymax=0.115,
  legend style={
    at={(0.5,-0.22)}, anchor=north,
    legend columns=2, font=\small,
    /tikz/every even column/.append style={column sep=0.5em}
  },
  grid=major, grid style={dotted, gray!50},
  tick label style={font=\small},
  label style={font=\small},
]
% 真値 θ₀=112 の縦線
\addplot[gray!60, dashed, thick] coordinates {(112,0) (112,0.115)};
\addlegendentry{真値 $\theta_0=112$}

% θ̂_e: Normal(112, 5.8138)
\addplot[black, thick, domain=90:135, samples=200]
  { gauss(112, 5.8138) };
\addlegendentry{$\hat\theta_e$（標本平均）}

% θ̂_b κ₀=1: Normal(110, 4.8447)
\addplot[blue!80!black, thick, domain=90:135, samples=200]
  { gauss(110, 4.8447) };
\addlegendentry{$\hat\theta_b,\ \kappa_0=1$}

% θ̂_b κ₀=2: Normal(108.571, 4.1526)
\addplot[red!80!black, thick, domain=90:135, samples=200]
  { gauss(108.571, 4.1526) };
\addlegendentry{$\hat\theta_b,\ \kappa_0=2$}

% θ̂_b κ₀=3: Normal(107.5, 3.6336)
\addplot[green!60!black, thick, domain=90:135, samples=200]
  { gauss(107.5, 3.6336) };
\addlegendentry{$\hat\theta_b,\ \kappa_0=3$}
\end{axis}
\end{tikzpicture}
\end{center}

\medskip
\noindent
\textbf{読み取れること：}
\begin{itemize}
  \item \textbf{図1}：$n$が小さいほどベイズ推定量の相対MSEが低い（事前情報が効く）．
        $n \to \infty$ でどの曲線も1に収束する（データが増えれば差がなくなる）．
  \item \textbf{図2}：$\kappa_0$が大きいほど分布が左（$\mu_0=100$寄り）にシフトし，
        かつ分布が狭くなる（分散減少）．$\hat\theta_e$は真値$\theta_0=112$に
        中心があるが分散が最大．
  \item バイアス（中心のズレ）と分散（幅）のトレードオフが視覚的に確認できる．
\end{itemize}

\end{exbox}

% ===========================
\section{まとめ}
% ===========================

\begin{thmbox}{5.4節のまとめ}
\begin{enumerate}
  \item \textbf{バイアス}は推定量の期待値と真値のズレ．不偏性は良い推定量の必要条件ではない．

  \item \textbf{MSE = 分散 + バイアス$^2$}というトレードオフがある．

  \item ベイズ推定量$\hat{\theta}_b$は標本平均$\hat{\theta}_e$よりバイアスがあるが分散が小さい．

  \item $(\mu_0 - \theta_0)^2 < \sigma^2(1/n + 2/\kappa_0)$が成立するとき，
  ベイズ推定量は MSE の意味で標本平均より優れる．

  \item 事前情報が多少でもあれば（$\mu_0$が大きく外れていなければ），
  特に$n$が小さい場合にベイズ推定量は有利になる．
\end{enumerate}
\end{thmbox}

\end{document}
